% PDF page 307
\source{gilbarg-trudinger:page-0307}

\setcounter{section}{11}
\noindent\textbf{12.1.\ Quasiconformal Mappings}

\begin{equation}\tag{12.2}
p_x^2+p_y^2+q_x^2+q_y^2\le 2K(p_xq_y-p_yq_x)+K',
\end{equation}

\noindent where $K$, $K'$ are constants, with $K\ge 1$, $K'\ge 0$. Although the geometric meaning
is no longer the same, we shall refer to the mappings obeying (12.2) as
$(K,K')$-quasiconformal. In the subsequent development it will be seen that
mappings satisfying (12.1) and (12.2) arise naturally from elliptic equations in
two variables, with $p$ and $-q$ representing the first derivatives of the solution.
The object of this section is the derivation of apriori interior H\"older estimates
for $(K,K')$-quasiconformal mappings. The main result will be the consequence of
lemmas concerning the Dirichlet integral

\begin{equation}\tag{12.3}
\D(r;z)=\iint_{B_r(z)} |Dw|^2\,dx\,dy=\iint_{B_r(z)}\bigl(|w_x|^2+|w_y|^2\bigr)\,dx\,dy
\end{equation}

\noindent of a $(K,K')$-quasiconformal mapping $w$ taken over disks $B_r(z)$. When there is no
ambiguity, we shall write $\D(r)$ for $\D(r;z)$ and $B_r$ for $B_r(z)$.

\medskip
\noindent\textbf{Lemma 12.1.} \emph{Let $w=p+iq$ be $(K,K')$-quasiconformal in a disk $B_R=B_R(z_0)$,
satisfying (12.2) with $K>1$, $K'\ge 0$, and suppose $|p|\le M$ in $B_R$. Then for all $r\le R/2$,}

\begin{equation}\tag{12.4}
\D(r)=\iint_{B_r}|Dw|^2\,dx\,dy\le C\left(\frac rR\right)^{2\alpha},
\qquad \alpha=K-(K^2-1)^{1/2},
\end{equation}

\noindent where $C=C_1(K)(M^2+K'R^2)$. If $K'=0$, the conclusion remains valid for $K=1$.

\medskip
\noindent\emph{Proof.} We first establish an estimate for the Dirichlet integral in the disk of radius
$R/2$. From (12.2) we have in any concentric disk $B_r\subset B_R$

\begin{equation}\tag{12.5}
\D(r)=\iint_{B_r}|Dw|^2\,dx\,dy\le 2K\iint_{B_r}\frac{\partial(p,q)}{\partial(x,y)}\,dx\,dy+K'\pi r^2
\end{equation}
\[
=2K\int_{C_r}p\,\frac{\partial q}{\partial s}\,ds+K'\pi r^2,
\]

\noindent where $s$ denotes arc length along the circle $C_r=\partial B_r$ described in the counter-
clockwise direction. Using the fact that $\D'(r)=\int_{C_r}|Dw|^2\,ds$, we observe

\begin{equation}\tag{12.6}
\int_{C_r}p\,\frac{\partial q}{\partial s}\,ds\le
\left(\int_{C_r}p^2\,ds\int_{C_r}|Dq|^2\,ds\right)^{1/2}
\end{equation}
\[
\le \left(\int_{C_r}p^2\,ds\int_{C_r}|Dw|^2\,ds\right)^{1/2}
\le M\bigl(2\pi r\,\D'(r)\bigr)^{1/2}.
\]


% PDF page 308
\source{gilbarg-trudinger:page-0308}



\noindent Inserting this into (12.5), and replacing $r$ by $R$ in the second member on the right,
we obtain

\medskip

\noindent (12.7)\qquad $\bigl[\D(r)-k_1\bigr]^2\leq k_2 r^2\D'(r),$

\medskip

\noindent where $k_1=\pi R^2K'$, $k_2=8\pi M^2K^2$. Now either $\D(R/2)\leq k_1$, in which case we have
the desired estimate; or if not, then $\D(r)>k_1$ for some $r=r_0<R/2$ and hence
for all larger $r$. The differential inequality (12.7) can then be integrated in
$r_0<r_1\leq r\leq r_2<R$ to yield

\[
\frac{1}{\D(r_1)-k_1}\geq \int_{r_1}^{r_2}\frac{\D(r)\,dr}{[\D(r)-k_1]^2}\geq \frac{1}{k_2}\log\frac{r_2}{r_1}.
\]

\noindent Taking $r_1=R/2$, $r_2=R$, we obtain

\medskip

\noindent (12.8)\qquad $\D(R/2)\leq \frac{8\pi}{\log 2}M^2K^2+\pi R^2K'.$

\medskip

\noindent We note that the derivation of this estimate involved no restrictions on $K, K'$
other than the non-negativity of $K$. We note also that it is not possible in general to
obtain such an estimate in the full disk $B_R$, as is shown by the set of analytic func-
tion $w_n=z^n$, $n=1,2,\dots,$ all of which satisfy $|w_n|\leq 1$ in $|z|\leq 1$, but

\[
\iint_{|z|<1}|Dw_n|^2\,dx\,dy\to\infty\qquad\text{as }n\to\infty;
\]

\noindent on the other hand, $\iint_{|z|<1-\delta}|Dw_n|^2\,dx\,dy\leq C(\delta)<\infty$ for any fixed $\delta>0$, where

\noindent $C(\delta)$ is independent of $n$.

\noindent We proceed now from the bound (12.8) on the Dirichlet integral in $B_{R/2}$ to a
growth estimate for $\D(r)$. From the inequalities,

\[
|p_xq_y|\leq \frac{\alpha}{2}p_x^2+\frac{1}{2\alpha}q_y^2,\qquad |p_yq_x|\leq \frac{\alpha}{2}q_x^2+\frac{1}{2\alpha}p_y^2,\qquad (\alpha>0),
\]

\noindent we obtain

\[
J=p_xq_y-p_yq_x\leq \frac{\alpha}{2}|w_x|^2+\frac{1}{2\alpha}|w_y|^2.
\]

\noindent Hence, writing (12.2) in the form,

\[
|w_x|^2+|w_y|^2\leq 2KJ+K',
\]


% PDF page 309
\source{gilbarg-trudinger:page-0309}

and substituting $\alpha = K - (K^2 - 1)^{1/2}$ (or, equivalently, $K = (1 + \alpha^2)/2\alpha$), we find

\[
|w_x|^2 \leq \frac{1}{\alpha^2}|w_y|^2 + \frac{2K'}{1 - \alpha^2}.
\]

Thus

\begin{equation}
|w_x|^2 = \frac{1}{1+\alpha^2}\bigl(|w_x|^2 + \alpha^2|w_x|^2\bigr)
\leq \frac{1}{1+\alpha^2}\left(|Dw|^2 + \frac{2\alpha^2K'}{1-\alpha^2}\right).
\tag{12.9}
\end{equation}

Since (12.2) is invariant under rotation, this inequality remains valid if any directional derivative $w_s$ replaces $w_x$.

We shall apply (12.9) to obtain a more precise estimation of $\int_{C_r} p q_s\, ds$ in (12.5).

Let $\bar p = \bar p(r)$ denote the mean value of $p$ over the circle $C_r$. Then

\begin{equation}
\int_{C_r} p q_s\, ds = \int_{C_r} (p-\bar p)q_s\, ds \leq \frac{1}{2}\int_{C_r}\left[\frac{(p-\bar p)^2}{r} + r q_s^2\right] ds.
\tag{12.10}
\end{equation}

We now make use of the Wirtinger inequality [HLP], which states

\[
\int_0^{2\pi}\bigl[p(r,\theta)-\bar p\bigr]^2\, d\theta \leq \int_0^{2\pi} p_\theta^2\, d\theta,
\]

that is,

\begin{equation}
\int_{C_r} (p-\bar p)^2\, ds \leq r^2 \int_{C_r} p_s^2\, ds.
\tag{12.11}
\end{equation}

(This result is easily proved by expanding $p = p(r,\theta)$ in a Fourier series in $\theta$ and applying Parseval’s equality.) Inserting (12.11) into (12.10), we see that

\[
\int_{C_r} p q_s\, ds \leq \frac{1}{2}\int_{C_r}(p_s^2 + q_s^2)\, ds = \frac{r}{2}\int_{C_r}|w_s|^2\, ds
\]

and hence by (12.9),

\[
\int_{C_r} p q_s\, ds \leq \frac{r}{2(1+\alpha^2)}\int_{C_r}|Dw|^2\, ds + \frac{2\pi\alpha^2K'}{1-\alpha^4}\,r^2.
\]


% PDF page 310
\source{gilbarg-trudinger:page-0310}

Now substituting this inequality into (12.5), and again using the relation

\[
\int_{C_r} |Dw|^2\, \dd s = \D'(r),
\]

we arrive at the differential inequality,

\begin{equation}\tag{12.12}
\D(r)\le \frac{r}{2\alpha}\,\D'(r)+kr^2,\qquad
k=\pi K'\left(1+\frac{2\alpha}{1-\alpha^2}\right).
\end{equation}

This implies

\[
-\frac{\dd}{\dd r}\bigl(r^{-2\alpha}\D(r)\bigr)\le 2\alpha k\, r^{1-2\alpha},
\]

from which it follows by integration between $r$ and $r_0$

\begin{equation}\tag{12.13}
\D(r)\le \left[\D(r_0)+\frac{\alpha}{1-\alpha}kr_0^2\right]
\left(\frac{r}{r_0}\right)^{2\alpha}.
\end{equation}

Letting $r_0=R/2$ and inserting the bound (12.8), we obtain the desired estimate
(12.4), with $C=\max\{C_2(K),C_3(K)\}(M^2+K'R^2)$, where

\[
C_2=\frac{32\pi}{\log 2}\,K^2,\qquad
C_3=\pi\left[4+\frac{\alpha}{1-\alpha}\left(1+\frac{2\alpha}{1-\alpha^2}\right)\right].
\]

We observe finally that when $K'=0$ the arguments are unaffected by allowing
$K=1$, and $C$ reduces to $C_2M^2$. $\square$

The following calculus lemma of Morrey provides the essential step from an
estimate of the growth of the Dirichlet integral to a Hölder estimate on the function
itself.

\medskip

\noindent\textbf{Lemma 12.2.} \textit{Let $w\in C^1(\Omega)$, and let $\widetilde{\Omega}\Subset\Omega$ with
$\dist(\widetilde{\Omega},\partial\Omega)>R$. Suppose there are positive constants $C$,
$\alpha$ and $R'$ such that}

\[
\D(r;z)=\iint_{B_r(z)} |Dw|^2\, \dd x\, \dd y\le Cr^{2\alpha}
\]

\textit{for all disks $B_r(z)$ with center $z\in \widetilde{\Omega}$ and radius
$r\le R'\le R$. Then for all $z_1,z_2\in \widetilde{\Omega}$ such that
$|z_2-z_1|\le R'$, we have}

\[
|w(z_2)-w(z_1)|\le 2\sqrt{\frac{C}{\alpha}}\,|z_2-z_1|^\alpha.
\]


% PDF page 311
\source{gilbarg-trudinger:page-0311}

This lemma is an immediate corollary of Theorem 7.19 for the case $n=2$ after applying the Schwarz inequality. A proof of the lemma in the form stated above is given in [FS].

The preceding lemmas are collected in the following apriori H\"older estimate for $(K, K')$-quasiconformal mappings.

\noindent\textbf{Theorem 12.3.} \textit{Let $w$ be $(K, K')$-quasiconformal in a domain $\Omega$, with $K>1$, $K'\geq 0$, and suppose $|w|\leq M$. Let $\tilde{\Omega}\subset \subset \Omega$ with $\operatorname{dist}(\tilde{\Omega}, \partial\Omega)>d$. Then for all $z_1, z_2\in \tilde{\Omega}$, we have}

\begin{equation}
|w(z_2)-w(z_1)|\leq C\left|\frac{z_2-z_1}{d}\right|^{\alpha}, \qquad \alpha=K-(K^2-1)^{1/2},
\tag{12.14}
\end{equation}

\noindent\textit{where $C=C_1(K)(M+d\sqrt{K'})$. If $K'=0$, then $C=C_1(K)M$ and the conclusion is valid also for $K=1$.}

\medskip

\noindent\textit{Proof.} Suppose first that $|z_2-z_1|\leq d/2$. The conditions of Lemmas 12.1 and 12.2 then apply with $R=d$ and $R'=d/2$, so that we have

\begin{equation*}
|w(z_2)-w(z_1)|\leq L\left|\frac{z_2-z_1}{d}\right|^{\alpha},
\end{equation*}

\noindent where

\[
L=C(K)(M^2+K'd^2)^{1/2}\leq C(K)(M+d\sqrt{K'}).
\]

If $|z_2-z_1|>d/2$, then

\begin{equation*}
|w(z_2)-w(z_1)|\leq 2M\leq 2M\left|\frac{z_2-z_1}{\tfrac12 d}\right|^{\alpha}<4M\left|\frac{z_2-z_1}{d}\right|^{\alpha}.
\end{equation*}

\noindent The theorem is therefore proved with $C_1(K)=\max(4, C(K))$. \hfill $\square$

\medskip

\noindent\textit{Remarks.} (1) The exponent $\alpha=K-(K^2-1)^{1/2}$ is the best (i.e. the largest) for which Lemma 12.1 and Theorem 12.3 are true. This can be seen from the example of the $K$-quasiconformal mapping $w(z)=r^{\alpha}e^{i\theta}$, $\alpha=K-(K^2-1)^{1/2}$, which has precisely the H\"older exponent $\alpha$ at $z=0$. The same results (for $K\geq 1$, $K'\geq 0$) with a smaller exponent $\alpha$ can be obtained from a slightly more direct proof of Lemma 12.1 starting with (12.5) and dispensing with (12.9); in this case the sharp form of the Wirtinger inequality (12.11) is not required (cf. [NI 1]).

\begin{enumerate}
\setcounter{enumi}{1}
\item Counterexamples show that Lemma 12.1 and Theorem 12.3 are not true for the exponent $\alpha=K-(K^2-1)^{1/2}$ when $K=1$ and $K'>0$, that is, for $\alpha=1$ (see Problem 12.1). However, if a mapping satisfies (12.2) with $K=1$, it satisfies such an inequality with any larger value of $K$ and the corresponding results in Lemma 12.1 and Theorem 12.3 then apply with exponent $\alpha$ arbitrarily close to $1$.
\end{enumerate}


% PDF page 312
\source{gilbarg-trudinger:page-0312}

\begin{quote}
(3) If $\tilde{\Omega}$ is bounded and can be covered by $N$ disks of diameter $d/2$, one sees from the proof that Theorem 12.3 remains valid under the weaker hypothesis
$
|p|\le M,
$
with the constant $C$ in (12.14) now depending also on $N$ and hence on the diameter of $\tilde{\Omega}$.

(4) \emph{Global estimates}. If $w=p+iq$ is $(K, K')$-quasiconformal in a $C^1$ domain $\Omega$ and $w\in C^1(\overline{\Omega})$, then Theorem 12.3 can be strengthened to an a priori global Hölder estimate for $w$. In particular, if $|w|\le M$ and $p=0$ on $\partial\Omega$, then $w$ satisfies a global Hölder condition with Hölder coefficient and exponent depending only on $K, K', M$ and $\Omega$. To outline the proof, let $\partial\Omega$ be the union of a finite number of overlapping arcs each of which can be straightened by a suitable $C^1$ diffeomorphism $(x,y)\to(\xi,\eta)$ defined in the neighbourhood of the arc. The function $w$ is quasiconformal in the $(\xi,\eta)$ variables with constants $\kappa, \kappa'$ depending on $K, K'$ and $\Omega$. By reflection across $\eta=0$, so that $p(\xi,-\eta)=-p(\xi,\eta)$ and $q(\xi,-\eta)=q(\xi,\eta)$ in the extended $(\xi,\eta)$ plane, the function $p+iq$ is seen to define a $(\kappa,\kappa')$-quasiconformal mapping to which the preceding interior estimates apply. Returning to the $(x,y)$ plane, we thus obtain a Hölder estimate for $w$ valid in $\overline{\Omega}$; that is
$$
|w(z_1)-w(z_2)|\le C|z_1-z_2|^\alpha,\qquad z_1,z_2\in\overline{\Omega},
$$
where $\alpha=\alpha(K,K',\Omega)$ and $C=C(K,K',\Omega,M)$. If $p=\tilde p$ on $\partial\Omega$, where $\tilde p\in C^1(\overline{\Omega})$ and $|\tilde p|_{1,\Omega}\le M'$, then by considering $p-\tilde p$ in place of $p$, we see that $w$ satisfies the same kind of global estimate, with $\alpha$ and $C$ now depending on $K,K',M,M'$, and $\Omega$.
\end{quote}

\section*{12.2. Hölder Gradient Estimates for Linear Equations}

The results of the preceding section will now be applied to obtain interior Hölder estimates for the first derivatives of solutions of uniformly elliptic equations,

\begin{equation}
\tag{12.15}
Lu = au_{xx}+2bu_{xy}+cu_{yy}=f,
\end{equation}

where $a$, $b$, $c$, $f$ are defined in a domain $\Omega$ of the $z=(x,y)$ plane. Let $\lambda=\lambda(z)$, $\Lambda=\Lambda(z)$ denote the eigenvalues of the coefficient matrix, so that

\begin{equation}
\tag{12.16}
\lambda(\xi^2+\eta^2)\le a\xi^2+2b\xi\eta+c\eta^2\le \Lambda(\xi^2+\eta^2)\qquad \forall(\xi,\eta)\in\mathbb{R}^2;
\end{equation}

and assuming $L$ is uniformly elliptic in $\Omega$, we have

\begin{equation}
\tag{12.17}
\frac{\Lambda}{\lambda}\le \gamma
\end{equation}

for some constant $\gamma\ge 1$. We suppose also that $\sup_{\Omega}(|f|/\lambda)\le \mu<\infty$. Dividing (12.15) by the minimum eigenvalue $\lambda$, we may assume that $\lambda=1$ and that (12.16) holds

% PDF page 313
\source{gilbarg-trudinger:page-0313}

with $\lambda=1$ and $\Lambda=\gamma$, while $|f|\le \mu$. Let us make this assumption in the following.

Setting

\[
p=u_x,\qquad q=u_y,
\]

we can write (12.15) as the system

\begin{equation}\tag{12.18}
\frac{a}{c}p_x+\frac{2b}{c}p_y+q_y=\frac{f}{c},\qquad p_y=q_x.
\end{equation}

By formal differentiation, $p$ is seen to be a solution (in the weak sense) of the uni-
formly elliptic equation of divergence form,

\[
\left(\frac{a}{c}p_x+\frac{2b}{c}p_y-\frac{f}{c}\right)_x+(p_y)_y=0,
\]

and a similar equation holds for $q$; (see the proof of Theorem 11.5). The H\"older
estimates on $p$ and $q$ that are derived in this section can also be obtained from the
methods developed in Chapter 8 for equations of divergence form. For $n=2$ the
details are not fundamentally different from those based on quasiconformal
mappings presented here.

Multiplying the left member of (12.18) by $cp_x$, we obtain

\[
p_x^2+p_y^2\le ap_x^2+2bp_xp_y+cp_y^2=cJ+fp_x,\qquad J=q_xp_y-q_yp_x,
\]

and similarly

\[
q_x^2+q_y^2\le aJ+fq_y.
\]

Adding these inequalities and noting that $2\le a+c=1+\Lambda\le 1+\gamma$, we have

\begin{equation}\tag{12.19}
|Dp|^2+|Dq|^2\le (a+c)J+f(p_x+q_y)
\le (1+\gamma)J+\frac12(1+\gamma)\mu(|p_x|+|q_y|).
\end{equation}

Inserting the inequality,

\[
(1+\gamma)\mu(|p_x|+|q_y|)\le \varepsilon(p_x^2+q_y^2)+\frac{1}{2\varepsilon}(1+\gamma)^2\mu^2,\qquad \varepsilon>0,
\]

and fixing $\varepsilon=1$ (the particular choice of $\varepsilon<2$ is inessential for our purposes),
we obtain from (12.19).

\begin{equation}\tag{12.20}
|Dp|^2+|Dq|^2\le 2(1+\gamma)J+(1+\gamma)^2\mu^2/2.
\end{equation}

% PDF page 314
\source{gilbarg-trudinger:page-0314}

Hence $w=p-iq$ (or $q+ip$) defines a $(K, K')$-quasiconformal mapping, satisfying
(12.2) with

\begin{equation}
\tag{12.21}
K = 1+\gamma, \qquad K' = (1+\gamma)^2\mu^2/2.
\end{equation}

By taking $\varepsilon$ sufficiently small, the constant $K$ can be made arbitrarily close to
$(1+\gamma)/2$.

If $f=0$, one obtains directly from (12.16) and (12.17) the inequality

$$
|Dw|^2 = |Dp|^2 + |Dp|^2 \le (a+c)J \le (1+\gamma)J.
$$

In this case the mapping $w=p-iq$ is $K$-quasiconformal with the constant
$K=(1+\gamma)/2$. An elementary but more careful calculation shows that the smallest
quasiconformality constant does not exceed $K=(\gamma+1/\gamma)/2$ (see Problem 12.3,
also [TA 1]).

We now establish the basic estimate for solutions of (12.15) required for the
ensuing nonlinear theory. We use the notation $d_z=\operatorname{dist}(z,\partial\Omega)$, $d_{1,2}=\min(d_{z_1},d_{z_2})$
and the interior norms and seminorms defined in (4.17) and (6.10); in particular,

$$
[u]_{1,\alpha}^*=\sup_{z_1,z_2\in\Omega} d_{1,2}^{1+\alpha}\frac{|Du(z_2)-Du(z_1)|}{|z_2-z_1|^\alpha}.
\qquad
\|f/\lambda\|_0^{(2)}=\sup_{z\in\Omega} d_z^2|f/\lambda|.
$$

\textbf{Theorem 12.4.} \emph{Let $u$ be a bounded $C^2(\Omega)$ solution of}

$$
Lu=au_{xx}+2bu_{xy}+cu_{yy}=f,
$$

\emph{where $L$ is uniformly elliptic, satisfying (12.16), (12.17) in a domain $\Omega$ of $\mathbb{R}^2$. Then
for some $\alpha=\alpha(\gamma)>0$, we have}

\begin{equation}
\tag{12.22}
[u]^*_{1,\alpha}\le C(\|u\|_0+\|f/\lambda\|_0^{(2)}), \qquad C=C(\gamma).
\end{equation}

The significant feature of this result is that the estimate (12.22) depends only
on bounds on the coefficients and not on any regularity properties. This is in
contrast with the Schauder estimates (Theorem 6.2) which depend as well on the
Hölder constants of the coefficients. The Hölder estimates of Chapter 8 for diver-
gence form equations in $n$ variables (Theorem 8.24) are also independent of regu-
larity properties of the coefficients, but those estimates concern the solution itself
and not its derivatives. The validity of the analogue of Theorem 12.4 for $n>2$
remains in doubt.

\emph{Proof of Theorem 12.4.} Let $z_1, z_2$ be any pair of points in $\Omega$, set $2d=d_{1,2}$, and
define $\Omega'=\{z\in\Omega\mid d_z>d\}, \Omega''=\{z\in\Omega' \mid \operatorname{dist}(z,\partial\Omega')>d\}$. We note that $z_1, z_2\in\bar{\Omega}''$.
We now apply Theorem 12.3 with $\Omega'$, $\Omega''$ in place of $\Omega$, $\tilde{\Omega}$ respectively and
$K=1+\gamma$, $K'=[(1+\gamma)\sup |f/\lambda|]^2/2$. The inequality (12.14) for $w=p-iq$, stated

% PDF page 315
\source{gilbarg-trudinger:page-0315}

in terms of the gradient $Du$, becomes for $\alpha = (1+\gamma) - (\gamma^2 + 2\gamma)^{1/2}$

\[
d^\alpha \frac{|Du(z_2) - Du(z_1)|}{|z_2 - z_1|^\alpha} \le C\bigl(\sup_{\Omega'} |Du| + d \sup_{\Omega'} |f/\lambda|\bigr), \qquad (C = C(\gamma)),
\]

\[
\le \frac{C}{d}\bigl(\sup_{\Omega'} d_z |Du(z)| + \sup_{\Omega'} d_z |f/\lambda|\bigr);
\]

hence

\[
d_{1,2}^{1+\alpha} \frac{|Du(z_2) - Du(z_1)|}{|z_2 - z_1|^\alpha} \le C\bigl(\sup_{\Omega} d_z |Du(z)| + \sup_{\Omega} d_z^2 |f/\lambda|\bigr),
\]

which implies

\[
[u]^*_{1,\alpha;\Omega'} \le C\bigl([u]^*_{1;\Omega} + |f/\lambda|_0^{(2)}\bigr), \qquad \text{for any } \Omega' \Subset \Omega.
\]

The interpolation inequality (6.8) for $j = k = 1$, $\beta = 0$, namely,

\begin{equation}
\tag{12.23}
[u]^*_{1} \le \varepsilon [u]^*_{1,\alpha} + C_1 |u|_0, \qquad (C_1 = C_1(\varepsilon)),
\end{equation}

gives

\[
[u]^*_{1,\alpha} \le C\bigl(\varepsilon [u]^*_{1,\alpha} + C_1 |u|_0 + |f/\lambda|_0^{(2)}\bigr).
\]

Choosing $\varepsilon$ so that $C\varepsilon = \frac{1}{2}$, we obtain for an appropriate constant $C = C(\gamma)$

\[
[u]^*_{1,\alpha} \le C\bigl(|u|_0 + |f/\lambda|_0^{(2)}\bigr),
\]

which gives the required result (12.22). \(\square\)

From (12.22) and the interpolation inequality (12.23) follows the norm estimate

\begin{equation}
\tag{12.24}
|u|^*_{1,\alpha} \le C\bigl(|u|_0 + |f/\lambda|_0^{(2)}\bigr), \qquad C = C(\gamma).
\end{equation}

\textbf{Global Estimates}

Theorem 12.4 can be extended to a $C^{1,\alpha}(\bar{\Omega})$ estimate under suitable smoothness hypotheses on the boundary data and the solution itself. For suppose, in addition to the hypotheses of Theorem 12.4, that $u \in C^2(\bar{\Omega})$, where $\Omega$ is a $C^2$ domain, and that $u = 0$ on $\partial \Omega$. Then we can assert a global bound: $|u|_{1,\alpha;\Omega} \le C$, where $\alpha = \alpha(\gamma,\Omega)$ and $C = C(\gamma,\Omega, |u|_0, |f/\lambda|_0)$. We outline the proof. As a normalization, we set $v = u/(1 + |Du|_0)$, so that $v$ satisfies $Lv = f/(1 + |Du|_0)$ and $|Dv| \le 1$. Let the boundary curve $\partial \Omega$ be covered by a finite number of overlapping arcs, each of which can be straightened into a segment of $\eta = 0$ by a suitable $C^2$ diffeomorphism $(x,y) \mapsto (\xi,\eta)$ defined in the neighborhood of the arc. As in the derivation of

% PDF page 316
\source{gilbarg-trudinger:page-0316}


(12.20), the mapping $p=p(\xi,\eta)$, $q=q(\xi,\eta)$, where $p=v_{\xi}$, $q=-v_{\eta}$, is $(\kappa,\kappa')$-quasiconformal in $(\xi,\eta)$ with constants $\kappa=\kappa(\gamma,\Omega)$, $\kappa'=\kappa'(\gamma,\Omega,|f|/\lambda_0)$ (we recall $|Dv|\leq 1$). Also, $p=0$ on $\eta=0$. The same argument as in Remark 4 at the end of Section 12.1 shows that $p$ and $q$, and hence $Dv$, satisfy a global Hölder condition in $\Omega$, in which the Hölder exponent $\alpha$ depends only on $\gamma$ and $\Omega$, and the Hölder coefficient depends also on $|f/\lambda_0|$. Thus,

\[
[u]_{1,x}\leq C(1+|Du|_0), \qquad C=C(\gamma,\Omega,|f/\lambda_0|),
\]

and by interpolation (see Lemma 6.35), we obtain a bound

\[
|u|_{1,\alpha;\Omega}\leq C(\gamma,\Omega,|u|_0,|f/\lambda_0|), \qquad \alpha=\alpha(\gamma,\Omega).
\]

If $\varphi\in C^2(\overline{\Omega})$ and $u=\varphi$ on $\partial\Omega$, then by considering $u-\varphi$ in place of $u$ in the preceding, and recalling that $|u|_0$ can be estimated in terms of sup$_{\partial\Omega}|\varphi|$ and $|f/\lambda_0|$ (Theorem 3.7), we infer the apriori global bound

\[
|u|_{1,\alpha;\Omega}\leq C=C(\gamma,\Omega,|\varphi|_2,|f/\lambda_0|), \qquad \alpha=\alpha(\gamma,\Omega).
\]

It should be emphasized that this estimate is independent of any regularity properties of $f$ and the coefficients of $L$. It is clear from the details that the dependence on $\Omega$ can be stated in terms of its dimensions and the bounds on the first and second derivatives of the mapping $\xi=\xi(x,y)$, $\eta=\eta(x,y)$, that is, in terms of the $C^2$ properties of $\partial\Omega$.

\noindent\hspace*{2em}Later in this chapter we make the following application of the above result.

Let $\Omega$ be a $C^{2,\beta}$ domain for some $\beta>0$ and let $f$ and the coefficients of $L$ lie in $C^\beta(\Omega)$. Suppose $u\in C^2(\Omega)\cap C^0(\overline{\Omega})$, $\varphi\in C^{2,\beta}(\overline{\Omega})$, satisfy $Lu=f$ in $\Omega$, $u=\varphi$ on $\partial\Omega$. Then $u\in C^{1,\alpha}(\overline{\Omega})$ and $|u|_{1,\alpha;\Omega}\leq C$, where $\alpha=\alpha(\gamma,\Omega)$ and $C=C(\gamma,\Omega,|\varphi|_2,|f/\lambda_0|)$. We note that $u$ is assumed only continuous on $\partial\Omega$. The proof follows from the preceding paragraph by an approximation argument. Namely, if $a_m$, $b_m$, $c_m$, $f_m$, $m=1,2,\dots$, are suitably chosen functions in $C^\beta(\overline{\Omega})$ converging to $a$, $b$, $c$, $f$, uniformly in compact subdomains of $\Omega$, the corresponding solutions $u_m$ of the Dirichlet problems, $L_m u_m=f_m$ in $\Omega$, $u_m=\varphi$ on $\partial\Omega$, are in $C^2(\overline{\Omega})$ and (by the preceding) satisfy a uniform $C^{1,\alpha}(\overline{\Omega})$ bound $|u_m|_{1,\alpha}\leq C$ for some $\alpha$ and $C$ independent of $m$. By the Schauder interior estimates and uniqueness, the sequence $\{u_m\}$ converges to the given solution $u$ of $Lu=f$. It follows that $u$ also satisfies the same $C^{1,\alpha}(\overline{\Omega})$ bound $|u|_{1,\alpha}\leq C$, which was our assertion.

\section*{12.3. The Dirichlet Problem for Uniformly Elliptic Equations}

In this section we prove existence by a procedure that is a variant of the one outlined in Chapter 11. In the Dirichlet problem treated here the details are generally simpler than in later problems and some steps of the program in Chapter 11 can be omitted.

We consider the Dirichlet problem for quasilinear elliptic equations of the


% PDF page 317
\source{gilbarg-trudinger:page-0317}

\begin{quote}
\textbf{general form,}

\begin{equation}\tag{12.25}
Qu = a(x,y,u,u_x,u_y)u_{xx} + 2b(x,y,u,u_x,u_y)u_{xy}
 + c(x,y,u,u_x,u_y)u_{yy} + f(x,y,u,u_x,u_y)=0,
\end{equation}

defined in a bounded domain $\Omega$ in the $(x,y)$ plane. Concerning the operator $Q$ we
shall assume:

\begin{enumerate}
\item[(i)] The functions $a=a(x,y,u,p,q), \dots, f=f(x,y,u,p,q)$ are defined for all
$(x,y,u,p,q)$ in $\Omega\times\mathbb{R}\times\mathbb{R}^2$ and, in addition, $a,b,c,f\in C^\beta(\Omega\times\mathbb{R}\times\mathbb{R}^2)$ for some
$\beta\in(0,1)$.

\item[(ii)] The operator $Q$ is uniformly elliptic in $\Omega$ for bounded $u$; that is, the
eigenvalues $\lambda=\lambda(x,y,u,p,q)$, $\Lambda=\Lambda(x,y,u,p,q)$ of the coefficient matrix satisfy
\begin{equation}\tag{12.26}
1\le \frac{\Lambda}{\lambda}\le \gamma(|u|)\qquad \forall (x,y,u,p,q)\in\Omega\times\mathbb{R}\times\mathbb{R}^2,
\end{equation}
where $\gamma$ is non-decreasing.

\item[(iii)] The function $f$ satisfies the structure conditions,
\begin{equation}\tag{12.27}
\frac{|f|}{\lambda}\le \mu(|u|)(1+|p|+|q|)
\end{equation}
\begin{equation}\tag{12.28}
\frac{f}{\lambda}\operatorname{sign}u\le \nu(1+|p|+|q|)\qquad \forall (x,y,u,p,q)\in\Omega\times\mathbb{R}\times\mathbb{R}^2,
\end{equation}
where $\mu$ is non-decreasing and $\nu$ is a non-negative constant. These correspond to
the conditions on lower order terms in linear equations. Equations satisfying
(12.28) are discussed in Chapter 10.
\end{enumerate}

We now establish the following existence theorem.

\textbf{Theorem 12.5.} \textit{Let $\Omega$ be a domain in $\mathbb{R}^2$ satisfying an exterior sphere condition,
and let $\varphi$ be a continuous function on $\partial\Omega$. Then if $Q$ is an elliptic quasilinear operator
satisfying conditions (i)--(iii), the Dirichlet problem}

\begin{equation}\tag{12.29}
Qu=0\text{ in }\Omega,\qquad u=\varphi\text{ on }\partial\Omega,
\end{equation}

\textit{has a solution $u\in C^{2,\beta}(\Omega)\cap C^0(\overline{\Omega})$.}

\textit{Proof.} We first prove the theorem under the more restrictive hypothesis,

\begin{equation}\tag{12.30}
\frac{|f|}{\lambda}\le \mu<\infty\qquad (\mu=\text{const}),
\end{equation}

in place of (12.27). The argument is based on reduction to the Schauder fixed
point theorem (Theorem 11.1). To define the mapping $T$ appearing in the state-
ment of that theorem, we make the following observation. Let $v$ be any bounded
\end{quote}

% PDF page 318
\source{gilbarg-trudinger:page-0318}

function with locally H\"older continuous first derivatives in $\Omega$, and let $\bar a=\bar a(x, y)=a(x, y, v, v_x, v_y), \ldots,$ be the locally H\"older continuous functions in $\Omega$ obtained by inserting $v$ for $u$ in the coefficients of $Q$. Since $|f|/\lambda$ is bounded, it follows from Theorem 6.13 that the linear Dirichlet problem,

\begin{equation}
\tag{12.31}
\bar a u_{xx}+2\bar b u_{xy}+\bar c u_{yy}+\bar f=0 \text{ in } \Omega,\qquad u=\varphi \text{ on } \partial\Omega,
\end{equation}

has a unique solution $u\in C^2(\Omega)\cap C^0(\bar\Omega)$. We observe from Theorem 3.7 that

\[
|u|_0=\sup_{\Omega}|u|\leq \sup_{\partial\Omega}|\varphi|+C_1\mu=M_0,\qquad C_1=C_1(\operatorname{diam}\Omega).
\]

Furthermore, if $\sup_{\Omega}|v|\leq M_0$ and if we set $\gamma_0=\gamma(M_0)$, we have from Theorem 12.4

\begin{equation}
\tag{12.32}
|u|_{1,\alpha}^*\leq C\bigl(|u|_0+\mu(\operatorname{diam}\Omega)^2\bigr),\qquad \alpha=\alpha(\gamma_0),\ C=C(\gamma_0),
\end{equation}
\[
\leq C\bigl(M_0+\mu(\operatorname{diam}\Omega)^2\bigr)=K.
\]

We note especially that this estimate depends only on the bound $M_0$ of the function $v$ used in defining the coefficients of equation (12.31).
Let us introduce the Banach space

\[
C_*^{1,\alpha}=C_*^{1,\alpha}(\Omega)=\{u\in C^{1,\alpha}(\Omega)\mid |u|_{1,\alpha;\Omega}^*<\infty\},
\]

where $\alpha$ is the H\"older exponent in (12.32). We can define a mapping $T$ of the set,

\[
\mathcal{G}=\{v\in C_*^{1,\alpha}\mid |v|_{1,\alpha}^*\leq K,\, |v|_0\leq M_0\},
\]

by letting $u=Tv$ be the unique solution of the linear Dirichlet problem (12.31) for $v\in\mathcal{G}$. By virtue of (12.32) and the bound $|u|_0\leq M_0$, we have $u\in\mathcal{G}$ and hence $T$ maps $\mathcal{G}$ into itself. Since $\mathcal{G}$ is convex and is closed in the Banach space

\[
C_*^1=\{u\in C^1(\Omega)\mid |u|_{1;\Omega}^*<\infty\},
\]

we may conclude from the Schauder fixed point theorem (Corollary 12.2) that $T$ has a fixed point, $u=Tu$, in $\mathcal{G}$ provided the mapping $T$ is continuous in $C_*^1$ and the image $T\mathcal{G}$ is precompact. This will provide a solution of the problem (12.29) under the hypothesis (12.30).
To prove $T\mathcal{G}$ is precompact in $C_*^1$, we observe first that the set $\mathcal{G}$, and hence $T\mathcal{G}$, is equicontinuous at each point of $\Omega$. We claim that the functions of $T\mathcal{G}$ are also equicontinuous at each point $z_0\in\partial\Omega$. For let $w$ be the barrier function in the argument preceding Theorem 6.13. The function $w$, which depends only on the ellipticity modulus $\gamma_0$ (in equation (12.31)) and the radius of the exterior disk at $z_0$, has the property that for any $\epsilon>0$ and a suitable constant $k_\epsilon$ independent of $v\in\mathcal{G}$,


% PDF page 319
\source{gilbarg-trudinger:page-0319}

the solution $u=Tv$ of (12.31) satisfies the inequality

\begin{equation}
\tag{12.33}
|u(z)-\varphi(z_0)|\leq \epsilon+k_{\epsilon}w(z) \qquad \text{in } \Omega.
\end{equation}

Since $w(z)\to 0$ as $z\to z_0$, this implies the equicontinuity of the set $T\mathcal{G}$ at $z_0$. The functions of $T\mathcal{G}$ are therefore equicontinuous on $\bar{\Omega}$. Since $T\mathcal{G}$ is a bounded equicontinuous set in $C_*^{1,\alpha}$, it is precompact in $C_*^1$ (see Lemma 6.33).

The continuity of $T$ in $C_*^1$ is proved in a similar way: Let $v,v_n\in\mathcal{G}$, $n=1,2,\ldots,$ and assume $|v_n-v|_1^*\to 0$ as $n\to\infty$. Consider $u=Tv$ and the sequence $u_n=Tv_n$, $n=1,2,\ldots$; we wish to show $|u_n-u|_1^*\to 0$. By the Schauder interior estimates and the Remark following Corollary 6.3, it follows that a suitable (renumbered) subsequence $\{u_m\}\subset\{u_n\}$ converges uniformly with its first and second derivatives in compact subsets of $\Omega$ to a solution $\tilde{u}$ in $\Omega$ of the limit equation (12.31) obtained by inserting $v$ into the coefficients of $Q$. We claim that $\tilde{u}(z)\to\varphi(z_0)$ for all $z_0\in\partial\Omega$ and hence (by uniqueness) $\tilde{u}=u=Tv$. For by the same barrier argument as above we may assert (12.33) with $u_m$ replacing $u$, from which we obtain in the limit
$\tilde{u}(z)\to\varphi(z_0)$ as $z\to z_0$. Thus $\tilde{u}=u$ and we have $Tu_m\to Tv$ on $\bar{\Omega}$ for the subsequence $\{v_m\}$.

Since the sequence $\{Tv_m\}$ is contained in $T\mathcal{G}$, which is precompact in $C_*^1$, a suitable subsequence of $\{Tv_m\}$ converges in the $C_*^1$ norm to $Tv$. The same argument, repeated on arbitrary subsequences of $\{v_n\}$, shows that $|Tv_n-Tv|_1^*\to 0$ for the entire sequence. This establishes the continuity of $T$ on $C_*^1$ and, as already observed, we may conclude the existence of a fixed point, $u=Tu$, in $\mathcal{G}$.

The theorem is thus proved in the special case that $|f|/\lambda$ is bounded on $\Omega\times\mathbb{R}\times\mathbb{R}^2$, in particular when $f\equiv 0$. We return now to the original hypothesis (iii). It will be convenient in the following to assume $\lambda(x,y,u,p,q)\equiv 1$ in $\Omega\times\mathbb{R}\times\mathbb{R}^2$, which can always be achieved by dividing the functions $a,b,c,f$ by $\lambda$. In this case (12.27), (12.28) become

\begin{equation}
\tag{12.34}
|f|\leq \mu(|u|)(1+|p|+|q|).
\end{equation}

\begin{equation}
\tag{12.35}
f\,\mathrm{sign}\,u\leq v(1+|p|+|q|),\qquad v=\mathrm{const}.
\end{equation}

We now proceed by truncation of $f$ to reduce the given problem (12.29) to the above case of bounded $f$. Namely, let $\psi_N$ denote the function given by

\[
\psi_N(t)=
\begin{cases}
t, & |t|\leq N\\
N\,\mathrm{sign}\,t, & |t|>N
\end{cases}
\]

and define the truncation of $f$ by

\[
f_N(x,y,u,p,q)=f\bigl(x,y,\psi_N(u),\psi_N(p),\psi_N(q)\bigr).
\]

From (12.34) we have $|f_N|\leq \mu N(1+2N)$. Consider now the family of problems,

\begin{equation}
\tag{12.36}
Q_Nu=a(x,y,u,Du)u_{xx}+2b(x,y,u,Du)u_{xy}+c(x,y,u,Du)u_{yy}
+f_N(x,y,u,Du)=0,
\end{equation}

\[
u=\varphi \qquad \text{on } \partial\Omega.
\]


% PDF page 320
\source{gilbarg-trudinger:page-0320}

\[
\text{308}
\qquad
\text{12.\ Equations in Two Variables}
\]

By virtue of (12.35) and Theorem 10.3, any solution $u$ in this family is subject to the
bound, independent of $N$,

\begin{equation}
\tag{12.37}
\sup_{\Omega} |u| \le \sup_{\partial \Omega} |\varphi| + C_1(v,\diam \Omega)=M.
\end{equation}

From the preceding treatment of problem (12.29) with bounded $f$ the problem
(12.36)---in which $f_N$ is bounded---is seen to have a solution
$u_N \in C^{1,\alpha}_*(\Omega)\cap C^{2,\beta}(\Omega)\cap C^0(\overline{\Omega})$,
$\alpha=\alpha(\gamma)$, $\gamma=\gamma(M)$. Furthermore, from Theorem 12.4 we infer
the estimate,

\begin{equation}
\tag{12.38}
[u_N]^*_{1,\alpha} \le C(|u_N|_0+|f_N|_0^{(2)}),
\end{equation}

where $C=C(\gamma)$, $f_N=f_N(x,y,u_N,Du_N)$. By (12.34) and (12.37), this becomes

\[
[u_N]^*_{1,\alpha} \le C(1+[u_N]^*_{1,\alpha}), \qquad C=C(M,\gamma,\mu,\diam \Omega), \ \mu=\mu(M).
\]

The interpolation inequality (12.23), with $C\varepsilon=\frac12$, now yields the uniform bound,
independent of $N$,

\begin{equation}
\tag{12.39}
|u_N|_{1,\alpha}^* \le C=C(M,\gamma,\mu,\diam \Omega).
\end{equation}

Applying the preceding estimate and the Schauder interior estimates (Corollary
6.3) to the family of equations $Q_nu_n=0$ on compact subsets, we obtain a (re-
numbered) subsequence $\{u_n\}\subset \{u_N\}$ that converges to a solution $u$ of $Qu=0$ in
$\Omega$ satisfying the estimate (12.39).

It remains to show that $u$ also satisfies the boundary condition $u=\varphi$ and for
this purpose we apply a barrier argument very similar to that given above. By virtue
of (12.34) and (12.37), each $u_n$ is the solution of a linear equation,

\[
Q_nu=a_n^{ij}D_{ij}u+b_n^iD_iu+f_n=0,\qquad i,j=1,2,
\]

where $a_n^{11}(x,y)=a(x,y,u_n,Du_n),\ldots,$ and where $b_n^i(x,y), f_n(x,y)$ are bounded
independently of $n$. At any point $z_0\in \partial \Omega$ the barrier argument preceding Theorem
6.13 applies to this family of equations with the barrier function $w$ in (6.45), which
depends only on $\gamma$, $\mu$ and the radius of the exterior disk at $z_0$. We therefore obtain
for any $\varepsilon>0$ and a suitable constant $k_\varepsilon$ independent of $n$ the inequality

\[
|u_n(z)-\varphi(z_0)|\le \varepsilon+k_\varepsilon w(z)\qquad \text{in } \Omega.
\]

Letting $n\to \infty$, we infer the same inequality with $u$ in place of $u_n$, and hence
$u(z)\to \varphi(z_0)$ as $z\to z_0$. This completes the proof of the theorem. $\square$

\medskip

\noindent\textit{Remarks.} \textit{(1) The proof of the preceding theorem is based only on interior
derivative estimates, thus allowing more general conditions on the coefficients and
boundary data. With the use of global estimates, a simple modification of the proof}

% PDF page 321
\source{gilbarg-trudinger:page-0321}

yields a $C^{2,\beta}(\bar{\Omega})$ solution when $\partial\Omega$ and $\varphi$ are in $C^{2,\beta}$ and $a,b,c,f \in C^\beta(\bar{\Omega}\times \mathbb{R}\times \mathbb{R}^2)$;
(cf. Problem 12.5). Under these same hypotheses the solution provided by Theorem 12.5 must also lie in $C^{2,\beta}(\bar{\Omega})$. To prove this assertion, we observe first that the equation $Qu=0$, after inserting $u$ into the coefficients, has coefficients in $C^\beta(\Omega)$, while $u \in C^{2,\beta}(\Omega)\cap C^0(\bar{\Omega})$ and $u=\varphi$ on $\partial\Omega$, where $\varphi\in C^{2,\beta}(\bar{\Omega})$. According to the results on global estimates for linear equations at the end of Section 12.2, the solution $u$ lies in $C^{1,\alpha}(\bar{\Omega})$ for some $\alpha$. The coefficients of $Qu$ are therefore in $C^{\alpha\beta}(\bar{\Omega})$. It follows from Theorem 6.14 and uniqueness that $u\in C^{2,\alpha\beta}(\bar{\Omega})$ and hence the coefficients of $Qu$ are in $C^\beta(\bar{\Omega})$. We infer in the same way that $u\in C^{2,\beta}(\bar{\Omega})$, as asserted.

\begin{enumerate}
\item[(2)] Condition (12.28) was imposed to insure a uniform bound on the magnitude of all possible solutions of the Dirichlet problem for $Q_Nu=0$. Whenever such a bound is known apriori, condition (12.28) may be omitted.

\item[(3)] The linear growth condition (12.27) on $|f|/\lambda$ was required to obtain the bound (12.39) by interpolation for $[u]_1^*$ in terms of $[u]_{1,\alpha}^*$ and $|u|_0$. Whenever an apriori bound on the gradient of solutions is known for the family of equations $Q_Nu=0$, this growth condition is superfluous; (such examples will be discussed in Chapter 15).

\item[(4)] The hypothesis that $\Omega$ satisfy an exterior sphere condition can be replaced by any other condition that guarantees the existence of a barrier for strictly elliptic linear equations $Lu=f$, where $f$ and the coefficients of $L$ are bounded. For example, it suffices that $\Omega$ satisfy an exterior cone condition (see Problem 6.3).
\end{enumerate}

\section*{12.4. Non-Uniformly Elliptic Equations}

In our study of non-uniformly elliptic equations we shall see that, unlike the preceding section in which the domain is essentially arbitrary, the solvability of the Dirichlet problem is in general closely connected with the geometry of the domain. This feature of non-uniformly elliptic problems has already been observed in the linear theory (Sect.\ 6.6). The results of the present section emphasize the important role of convexity of the domain in insuring solvability of the Dirichlet problem for general quasilinear elliptic equations of the form,

\begin{equation}
\tag{12.40}
Qu=a(x,y,u,u_x,u_y)u_{xx}+2b(x,y,u,u_x,u_y)u_{xy}+c(x,u,y,u,u_x,u_y)u_{yy}
=0.
\end{equation}

Let $\Omega$ be a bounded domain in $\mathbb{R}^2$ and $\varphi$ be a function defined on $\partial\Omega$. The Dirichlet problem for equation (12.40) will be formulated in terms of the boundary curve

\[
\Gamma=(\partial\Omega,\varphi)=\{(z,\varphi(z))\in \mathbb{R}^3\mid z\in \partial\Omega\}.
\]

We recall (as in Sect.\ 11.3) that $\Gamma$ and $\varphi$ satisfy a bounded slope condition (with


% PDF page 322
\source{gilbarg-trudinger:page-0322}

constant $K$) if for every point $P=(z_0,\varphi(z_0))\in \Gamma$ there are planes in $\mathbb{R}^3$,
\[
u=\pi_p^{\pm}(z)=a^{\pm}\cdot(z-z_0)+\varphi(z_0), \qquad a^{\pm}=a^{\pm}(z_0),
\]
passing through $P$ such that:
\begin{equation}
\begin{aligned}
\text{(i)}\quad &\pi_p^{-}(z)\leq \varphi(z)\leq \pi_p^{+}(z) \qquad \forall z\in \partial\Omega;\\
\text{(ii)}\quad &|D\pi_p^{\pm}|=|a^{\pm}(z_0)|\leq K \qquad \forall z_0\in \partial\Omega.
\end{aligned}
\tag{12.41}
\end{equation}

Condition (i) states that for each $P$ the curve $\Gamma$ is bounded above and below on the
cylinder $\partial\Omega\times \mathbb{R}$ by the planes $u=\pi_p^{+}(z)$ and $u=\pi_p^{-}(z)$, and coincides with them
at $P$. Condition (ii) states that the slopes of the planes are uniformly bounded,
independently of $P$, by the constant $K$. It is evident that the bounded slope condition
implies the continuity of $\varphi$.

We make the following remarks concerning the bounded slope condition.

\begin{enumerate}
\item Whatever the domain $\Omega$, $\Gamma$ lies in a plane (and satisfies a bounded slope
condition) if and only if $\varphi$ is the restriction of a linear function on $\partial\Omega$. However, if
$\Gamma$ does not lie in a plane and satisfies a bounded slope condition, then $\Omega$ must be
convex. For we have from (12.41)
\[
0\neq \pi_p^{+}(z)-\pi_p^{-}(z)=(a^{+}-a^{-})\cdot (z-z_0)\geq 0 \qquad \forall z\in \partial\Omega,
\]
and thus there is a supporting line $(a^{+}-a^{-})\cdot (z-z_0)=0$ at each $z_0\in \partial\Omega$, which
implies the convexity of $\Omega$.

\item Suppose $\partial\Omega$ is convex and $\Gamma=(\partial\Omega,\varphi)$ satisfies a bounded slope condition.
Let $P_i=(z_i,\varphi(z_i))$, $i=1,2,3$, be three distinct points on $\Gamma$. If $z_1$, $z_2$, $z_3$ are collinear
then $P_1$, $P_2$, $P_3$ are also collinear on $\Gamma$, for otherwise these points would determine
a vertical plane, contradicting (12.41). Thus $\varphi$ is linear on the straight segments of $\Omega$.

\item Closely related to and, in fact, equivalent to the bounded slope condition
is the following \emph{three-point condition}, which appears often in the literature on
minimal surfaces and non-parametric variational problems in two independent
variables. Let $\Omega$ be bounded and convex; then the curve $\Gamma=(\partial\Omega,\varphi)$ is said to
satisfy a three-point condition with constant $K$ if every set of three distinct points
on $\Gamma$ lies in a plane of slope $\leq K$. At the end of this section we prove the equivalence
of the bounded slope and three-point conditions, with the same constant $K$.

It follows from the three-point condition that the plane determined by any
three non-collinear points of $\Gamma$ must have slope $\leq K$. Thus, if $\Omega$ is strictly convex
(that is, the open straight segment joining any two points of $\partial\Omega$ lies entirely in $\Omega$),
then the slope of every plane intersecting $\Gamma$ in at least three points does not exceed $K$
and, conversely, this stronger form of the three-point condition obviously implies
that $\partial\Omega$ is strictly convex.

\item It is not difficult to show that if $\varphi\in C^2$, $\partial\Omega\in C^2$ and the curvature of $\partial\Omega$
is everywhere positive, then $\Gamma=(\partial\Omega,\varphi)$ satisfies a bounded slope condition, with
a constant depending on the minimum curvature of $\partial\Omega$ and bounds on the first
and second derivatives of $\varphi$; (see [SC 3], [HA]).
\end{enumerate}

% PDF page 323
\source{gilbarg-trudinger:page-0323}

The solution of the Dirichlet problem for (12.40) will require an apriori bound
for the gradient which is provided by the following lemma.

\paragraph{Lemma 12.6.}
Let $\Omega$ be a bounded domain in $\mathbb{R}^2$, and let $\varphi$ be a function defined on
$\partial\Omega$ satisfying a bounded slope condition with constant $K$. Suppose $u\in C^2(\Omega)\cap C^0(\bar{\Omega})$
satisfies a linear elliptic equation,

\begin{equation}
\tag{12.42}
Lu = au_{xx} + 2bu_{xy} + cu_{yy} = 0 \qquad \text{in } \Omega,
\end{equation}

with $u=\varphi$ on $\partial\Omega$. Then

\begin{equation}
\tag{12.43}
\sup_{\Omega} |Du| \le K
\end{equation}

We emphasize that $L$ is only required to be elliptic and that no other conditions
are placed on the coefficients. It is of interest that the result is valid even more
generally, for arbitrary saddle surfaces $u=u(x,y)$ (see [RA], [NU]).

\paragraph{Proof of Lemma 12.6.}
We observe first that if $\varphi$ is the restriction of a linear function
on $\partial\Omega$, then by uniqueness (Theorem 3.3) the solution $u$ coincides with this function
in $\Omega$ and the conclusion (12.43) holds trivially. It follows from the remark (1)
above that $\Omega$ may be assumed convex.

From (12.42) and the ellipticity of $L$ we have

\[
0\le au_{xx}^2+2bu_{xx}u_{xy}+cu_{xy}^2=c(u_{xy}^2-u_{xx}u_{yy}),
\]
\[
0\le au_{xy}^2+2bu_{xy}u_{yy}+cu_{yy}^2=a(u_{xy}^2-u_{xx}u_{yy}).
\]

Accordingly, $u_{xx}u_{yy}-u_{xy}^2\le 0$ and the equality holds only at points where $D^2u=0$.
(We note that $u=u(x,y)$ is a saddle surface.) Consider now any point $z_0=
(x_0,y_0)\in \Omega$ where $u_{xx}u_{yy}-u_{xy}^2<0$ and let $u_0(x,y)=Ax+By+C$ define the tangent
plane $\Pi$ to the surface $u=u(x,y)$ at $(x_0,y_0)$. The function $w=u-u_0$ is a solution
of (12.42) in $\Omega$, and the set on which $w=0$ divides a small disk about $z_0$ into
precisely four domains $D_1,\dots,D_4$ in which $w$ is alternately positive and negative,
say $w>0$ in $D_1,D_3$ and $w<0$ in $D_2,D_4$. Let $D_1'\supset D_1$, $D_3'\supset D_3$ be components of
the set in $\Omega$ where $w>0$, and let $D_2'\supset D_2$, $D_4'\supset D_4$ be similarly defined. Then
each of the domains $D_1',\dots,D_4'$ has at least two boundary points on $\partial\Omega$ where
$w=0$, for otherwise the weak maximum principle (Theorem 3.1) would imply $w\equiv 0$
in the domain. It follows that $\Pi$ intersects the boundary curve $\Gamma=(\partial\Omega,\varphi)$ in at
least four points. From remark (3) above we infer that if any three of these points
are not collinear, then $\Pi$ has slope not exceeding $K$. On the other hand, if the points
on $\partial\Omega$ where $w=0$ form a collinear set $\Sigma$, then (by remark (2)) $\varphi$ and $u$ are linear
and identical with $u_0$ on the straight segment of $\partial\Omega$ containing $\Sigma$. This would imply
$w\equiv 0$ in one of the domains $D_i'$, which is a contradiction. Thus the slope of the
tangent plane to the solution surface cannot exceed $K$ at points where $D^2u\neq 0$.

Now consider the set $S$ on which $D^2u=0$. We may assume $S\neq \Omega$, for otherwise
$u(x,y)$ would be linear and the conclusion is trivial. If $z_0\in S$ and $z_0$ is the limit


% PDF page 324
\source{gilbarg-trudinger:page-0324}

of points where $D^2u\neq 0$, then by continuity the tangent plane at $z_0$ must have
inclination not in excess of $K$. The remaining possibility is that $z_0$ is an interior
point of $S$, in which case $z_0$ is contained in an open component $G$ of the set in
which $D^2u=0$. On $G$, we have that $u$ is linear and that the tangent plane coincides
with the surface $u=u(x,y)$. Since any boundary point of $G$ that is interior to $\Omega$ is
the limit of points where $D^2u\neq 0$, we conclude $|Du|\le K$ everywhere on $G$ and in
particular at $z_0$. This completes the proof. $\square$

We are now in the position to establish the following existence theorem for
equation (12.40), which extends Theorem 11.5 for the case $n=2$.

{\bf Theorem 12.7.} \emph{Let $\Omega$ be a bounded domain in $\mathbb{R}^2$ and assume that the equation,}

\[
Qu=a(x,y,u,u_x,u_y)u_{xx}+2b(x,y,u,u_x,u_y)u_{xy}+c(x,y,u,u_x,u_y)u_{yy}
\]
\[
=0,
\]

\emph{is elliptic in $\Omega$ with coefficients $a, b, c \in C^\beta(\Omega\times \mathbb{R}\times \mathbb{R}^2)$ for some $\beta \in (0,1)$. Let $\varphi$
be a function defined on $\partial\Omega$ satisfying a bounded slope condition with constant $K$.
Then the Dirichlet problem}

\[
Qu=0 \text{ in } \Omega,\qquad u=\varphi \text{ on } \partial\Omega
\]

\emph{has a solution $u\in C^{2,\beta}(\Omega)\cap C^0(\overline{\Omega})$, with $\sup_\Omega |Du|\le K$.}

{\bf Proof.} It suffices to assume $\Omega$ is convex, since otherwise $\varphi$ is the restriction of a
linear function and the result is trivial. We divide the coefficients $a, b, c$ by the
maximum eigenvalue $\Lambda=\Lambda(x,y,u,p,q)$ and denote again by $Qu=0$ the equation
thus obtained. The operator $Q$ now has maximum eigenvalue $1$, but the minimum
eigenvalue $\lambda$ may approach zero in $\Omega\times \mathbb{R}\times \mathbb{R}^2$. We consider the family of equations

\[
(12.44)\qquad Q_\epsilon u=Qu+\epsilon \,\Delta u=0,\qquad \epsilon>0,
\]

with the boundary condition $u=\varphi$ on $\partial\Omega$. For each $\epsilon$ this equation is uniformly
elliptic and Theorem 12.5 asserts the existence of a solution $u_\epsilon\in C^{2,\beta}(\Omega)\cap C^0(\overline{\Omega})$
such that $u_\epsilon=\varphi$ on $\partial\Omega$. (Only the first part of the proof of Theorem 12.5 is needed
here.) By Lemma 12.6, the solution $u_\epsilon$ satisfies the uniform gradient estimate,

\[
(12.45)\qquad \sup_\Omega |Du_\epsilon|\le K,
\]

independently of $\epsilon$. From the maximum principle we also have $|u_\epsilon|\le \sup_{\partial\Omega} |\varphi|$.
As a consequence, in every subdomain $\Omega' \Subset \Omega$, the linear equations,

\[
(12.46)\qquad Q_\epsilon u=a_\epsilon u_{xx}+2b_\epsilon u_{xy}+c_\epsilon u_{yy}+\epsilon\,\Delta u=0
\]

% PDF page 325
\source{gilbarg-trudinger:page-0325}

12.4. Non-Uniformly Elliptic Equations \hfill 313

obtained by setting $a_{\epsilon}(x,y)=a(x,y,u_{\epsilon},D u_{\epsilon}),\dots,$ have minimum eigenvalues,
$\lambda_{\epsilon}=\lambda(x,y,y_{\epsilon},D u_{\epsilon})$ for $(x,y)\in\Omega'$, that are uniformly bounded below by a constant
$\lambda(\Omega')>0$ depending only on $\Omega'$; the upper eigenvalues are of course bounded by
$2$ for all $\epsilon<1$. Theorem 12.4 now asserts that in subsets $\Omega''\subset\subset\Omega'$ the solutions
of (12.46) with $u=\varphi$ on $\partial\Omega$, in particular the solutions $u_{\epsilon}$, satisfy a uniform H\"older
bound for the gradient (independent of $\epsilon$):

\[
[D u]_{x;\Omega''}\leq C,\qquad \alpha=\alpha(\lambda(\Omega')).
\]

where $C=C(\lambda(\Omega'),\sup_{\partial\Omega}|\varphi|,\operatorname{dist}(\Omega'',\partial\Omega'))$. Accordingly, the coefficients $a_{\epsilon}, b_{\epsilon}, c_{\epsilon}$
are locally H\"older continuous with exponent $\alpha\beta$ in $\Omega'$ and have uniform bounds
in $C^{\alpha\beta}(\overline{\Omega'})$. Since $\Omega'$ and $\Omega''$ are arbitrary it follows from Corollary 6.3 and the
accompanying Remark that the family of solutions $u_{\epsilon}$ of (12.46) are equicontinuous
with their first and second derivatives on compact subsets of $\Omega$ and hence, by the
usual diagonalization process, there is a sequence $\{u_{\epsilon_n}\}$ of the family $\{u_{\epsilon}\}$ con-
verging in $\Omega$ to a solution $u_0$ of $\Q u=0$ as $\epsilon_n\to 0$. The uniform gradient bound
(12.45) guarantees that the convergence is uniform on $\overline{\Omega}$ and hence $u_0=\varphi$ on $\partial\Omega$.
This completes the proof of the theorem. $\square$

\textit{Remarks.} \ (1) That some geometric condition on $\Omega$, such as convexity, must be
imposed in general is indicated by classical counterexamples such as that of the
minimal surface equation,

\[
(12.47)\qquad (1+u_y^2)u_{xx}-2u_xu_yu_{xy}+(1+u_x^2)u_{yy}=0,
\]

in the annulus, $a<r<b$, $r=(x^2+y^2)^{1/2}$. If the boundary condition is $\varphi=h$
$(=\mathrm{const}>0)$ on $r=a$, $\varphi=0$ on $r=b$, and $h$ is sufficiently small, the boundary
value problem has the well-known catenoid solution. However, if $h$ is sufficiently
large, there is no solution taking on the prescribed boundary values. It will be seen
in Chapter 14 that the Dirichlet problem for the minimal surface equation (12.47)
is solvable for arbitrary $C^2$ boundary values if and only if $\Omega$ is convex.
(2) The smoothness hypotheses implicit in the boundary slope condition cannot
be relaxed in general to allow continuous boundary values $\varphi$. Counterexamples
show that the Dirichlet problem need not have a solution for continuous boundary
values even when the boundary curve is a circle and the coefficients of the equation
are arbitrarily smooth; (see [FN 2]). Solvability for continuous boundary values
will be discussed in Chapters 15 and 16.
(3) The essential step in the proof of Theorem 12.7 is the reduction to the uni-
formly elliptic case, which is made possible by the existence of an apriori global
gradient bound (Lemma 12.6). Such gradient bounds can also be established under
suitable structure conditions on the operator $\Q$ and geometric assumptions con-
cerning the domain $\Omega$. They are discussed in Chapters 14 and 15. In the case of
convex $\Omega$ it is sometimes possible to replace the planes of the bounded slope
condition by suitable super- and subfunctions with respect to $\Q$ and $\varphi$ and to
proceed with the argument in essentially the same way. Thus, if $\Omega$ is convex and


% PDF page 326
\source{gilbarg-trudinger:page-0326}

the operator $Q$ is such that $\Lambda/\lambda \le \mu(|u|)(1 + p^2 + q^2)$, the Dirichlet problem for
equation (12.40) is solvable for arbitrary $\varphi \in C^2$ whether or not the bounded
slope condition is satisfied (see Section 14.2). This includes, for example, the
minimal surface equation (12.47).

\medskip

\noindent (4) If $Q$, $\partial\Omega$ and $\varphi$ satisfy the additional smoothness hypotheses of Theorem
11.4, that is, $a$, $b$, $c \in C^\beta(\overline{\Omega} \times \mathbb{R} \times \mathbb{R}^2)$, $\partial\Omega \in C^{2,\beta}$, $\varphi \in C^{2,\beta}(\partial\Omega)$, then the solution
provided by Theorem 12.7 lies in $C^{2,\beta}(\overline{\Omega})$. The argument proceeds essentially
as in Remark 1 following Theorem 12.5, after the observation that the gradient
bound $|Du| \le K$ makes the equation $Qu=0$ uniformly elliptic.

\bigskip

\noindent\textbf{Equivalence of the Bounded Slope and Three-Point Conditions}

\medskip

Assume first that $\Gamma=(\partial\Omega, \varphi)$ satisfies a bounded slope condition with constant
$K$, $\Omega$ being convex. Let $z_1$, $z_2$, $z_3$ be three non-collinear points on $\partial\Omega$, and let

\[
u=\pi_i^+(z)=a_i^+ \cdot (z-z_i)+\varphi(z_i), \qquad i=1, 2, 3,
\]

be the corresponding upper and lower planes at the points $P_i=(z_i, \varphi(z_i)) \in \Gamma$
satisfying (12.41). Also let

\[
u=\pi(z)=a\cdot z+b=a\cdot (z-z_i)+\varphi(z_i), \qquad i=1, 2, 3,
\]

be the plane passing through $P_1$, $P_2$, $P_3$. We wish to show $|a| \le K$. (If $z_1$, $z_2$, $z_3$,
and hence $P_1$, $P_2$, $P_3$, are collinear, the plane $u=\pi(z)$ with slope $\le K$ can be determined by continuity.) From (12.41) it follows

\begin{equation}
\tag{12.48}
a_i^- \cdot (z_j-z_i)\le a\cdot (z_j-z_i)\le a_i^+ \cdot (z_j-z_i), \qquad i,j=1, 2, 3.
\end{equation}

Since $z_1$, $z_2$, $z_3$ are the vertices of a non-degenerate triangle and $a$ is a vector in
$\mathbb{R}^2$, we have for some $i=1$, $2$, or $3$, either

\[
a=\pm \sum_j c_j(z_j-z_i), \qquad c_j \ge 0.
\]

From (12.48) we infer either

\[
|a|^2 \le a_i^+ \cdot \sum_j c_j(z_j-z_i)=a_i^+ \cdot a \le K|a|
\]

or

\[
|a|^2 \le a_i^- \cdot \sum_j c_j(z_i-z_j)=a_i^- \cdot a \le K|a|,
\]

and hence $|a|\le K$. Thus the bounded slope condition implies the three-point condition
with the same constant $K$.


% PDF page 327
\source{gilbarg-trudinger:page-0327}

Conversely, let $\Gamma = (\partial \Omega, \varphi)$ satisfy the three-point
condition with constant $K$ over the convex domain $\Omega$. Let
$A = (z_A, \varphi(z_A))$ be any point of $\Gamma$ such that $z_A$ is not an
interior point of a straight segment of $\partial \Omega$. There exists a
sequence of triangles with vertices at non-collinear points $A$, $B_i$, $C_i \in
\Gamma$, $i = 1, 2, \ldots$, such that $B_i$, $C_i \to A$ as $i \to \infty$
and the planes $\Delta_i$ determined by $A$, $B_i$, $C_i$ converge to a
limiting plane $\Delta$. It can be assumed that the segments $AB_i$ have a
limiting direction, which determines a straight line $L$ through $A$ lying in
$\Delta$ whose projection $L_z$ on the $z$ plane is a supporting line of
$\Omega$. Obviously $L$ coincides with the tangent to $\Gamma$ at $A$
whenever the tangent exists. Let $Q \in \Gamma$, $Q \notin L$, and let $\Pi$
denote the plane determined by $Q$ and $L$. The slope of $\Pi$ does not exceed
$K$ since the sequence of planes determined by $A$, $Q$ and $B_i$ have the same
limit slope as $\Pi$. Consider now the set of planes containing $L$ and the
points $Q \in \Gamma$, $Q \notin L$, and let these planes be defined by the
linear functions $u = \pi_Q(z)$. In the $z$ plane let $H$ denote the half-plane
on the side of $L$ containing $\Omega$. If $\pi_Q(z) \ge q(z)$ for some
$z \in H$, then the same inequality holds for all $z \in H$ and, in
particular, for all $z \in \partial \Omega$. Thus, there are upper and lower
planes at $A$ defined by

\[
\pi^+(z) = \sup_{\substack{Q \in \Gamma \\ Q \notin L}} \pi_Q(z), \qquad
\pi^-(z) = \inf_{\substack{Q \in \Gamma \\ Q \notin L}} \pi_Q(z), \qquad z \in H.
\]

The planes $u = \pi^+(z)$ and $u = \pi^-(z)$ by the preceding discussion have
slopes not exceeding $K$ and lie respectively above and below $\Gamma$ (over
$\partial \Omega$). Hence they satisfy the bounded slope condition at $A$ with
constant $K$. It is clear that if $A$ lies on a straight segment of $\Gamma$,
then the line containing this segment can replace $L$ in the above argument.
Thus the three-point condition implies the bounded slope condition with the
same constant $K$.

Notes

Hölder estimates for quasiconformal mappings (Section 12.1) have been derived
in various ways starting with Morrey [MY 1]; (for references see [FS]). The
development here is due to Finn and Serrin [FS], who derived the estimate
(12.14) with its optimal Hölder exponent and, when $K' = 0$ in (12.2), also
obtained this result with a Hölder coefficient $CM$ in which $C$ is an absolute
constant.

The basic $C^{1,\alpha}$ Hölder estimates for solutions of linear equations in
Section 12.2 are due to Morrey [MY 1] and, in simpler form, to Nirenberg
[NI 1]. The presentation here is a variant of the latter, but follows Morrey in
using growth estimates for the Dirichlet integral. The idea of applying these
a priori $C^{1,\alpha}$ estimates for linear equations to the Dirichlet problem
for the quasilinear equation (12.40) appears in [MY 1], which contains a flaw,
however. The idea was carried to completion and simplified in the details by
Nirenberg who applied the Schauder fixed point theorem in arriving at the
general existence theorem described below. This approach is the basis for the
development in Sections 12.3 and 12.4.

Until the late 1950's the theory of the Dirichlet problem for nonlinear elliptic
equations was confined largely to the case of two independent variables. The

% PDF page 328
\source{gilbarg-trudinger:page-0328}




\noindent pioneering work of Bernstein [BE 1--4], which was somewhat restrictive in its hypotheses, was extended by Leray and Schauder [SC 3], [LS], who obtained solutions of the Dirichlet problem for (12.40) under the assumption of $C^{2,\beta}$ coefficients and sufficiently smooth boundary data by using methods based (as in Bernstein's work) on a priori estimates of the second derivatives of solutions. These results were improved and simplified by the above contributions of Morrey and Nirenberg, the latter proving that if (12.40) is elliptic in $\overline{\Omega}$ with coefficients in $C^{\beta}(\overline{\Omega}\times \R \times \R^{2})$, where $\partial\Omega$ is a $C^{2,\beta}$ uniformly convex curve, and $\varphi\in C^{2,\beta}(\partial\Omega)$, then the Dirichlet problem for (12.40) has a solution in $C^{2,\beta}(\overline{\Omega})$; (if $\varphi\in C^{1,1}(\partial\Omega)$ then a solution exists in $C^{2,\beta}(\Omega)\cap C^{0}(\overline{\Omega})$). Theorem 12.7 extends this result by weakening the hypotheses on the coefficients and on the boundary data by requiring the latter to satisfy only a bounded slope condition, without additional regularity assumptions. The solution thus obtained also lies in $C^{2,\beta}(\overline{\Omega})$ when the hypotheses are the same as in Nirenberg's theorem (see Remark 4 after Theorem 12.7). The proof of Theorem 12.7 depends only on apriori interior $C^{1,\alpha}$ estimates for linear equations (Theorem 12.4). That such estimates suffice for an existence proof was observed already by Nirenberg ([NI 1], p. 146).

\noindent Concerning the Dirichlet problem for the uniformly elliptic equation (12.25), we mention the contributions of Bers and Nirenberg [BN], Ladyzhenskaya and Ural'tseva [LU 4] and von Wahl [WA]. In [BN], existence results--for both the Dirichlet and Neumann problems--are obtained in the context of $W^{2,2}$ solutions (which are also in $C^{2,\alpha}$ if the coefficients in (12.25) are in $C^{\alpha}$). It assumes a linear growth condition for $f(x,y,u,p,q)$ analogous to (12.27), and the methods of proof in [BN] and Theorem 12.5 are similar in being based on $C^{1,\alpha}$ estimates for linear equations and the Schauder fixed point theorem. Both [BN] and Theorem 12.5 make no apriori assumptions concerning the solutions. However, if a uniform bound on $\sup_{\Omega}|u|$ is assumed for all solutions $u$ of (12.29) under the following hypotheses: $a$, $b$, $c$, $f\in C^{\beta}(\overline{\Omega}\times \R\times \R^{2})$ in (12.25), $\partial\Omega\in C^{2,\beta}$, $\varphi\in C^{2,\beta}(\overline{\Omega})$, and $|f/\lambda|\leq C(1+p^{2}+q^{2})$ (weakening the condition (12.27)), then [WA] establishes an apriori $C^{1,\alpha}(\overline{\Omega})$ bound on solutions of the Dirichlet problem, thereby extending a similar result in [LU 4]. (This estimate does not follow from the methods of the present chapter without an apriori gradient bound.) Existence of a $C^{2,\beta}(\overline{\Omega})$ solution of the Dirichlet problem (12.29) can now be obtained by application of the Leray--Schauder fixed point theorem.

\noindent The proof of the gradient bound in Lemma 12.6 is suggested by the introductory remarks of Finn in [FN 2]. It is usually inferred from a theorem of Rad\'o [RA], which asserts that if $u=u(x,y)$ is a saddle surface over a convex domain $\Omega$ and $u$ is continuous on $\overline{\Omega}$ with boundary values satisfying a three-point condition with constant $K$, then $u$ satisfies a Lipschitz condition with constant $K$ in $\Omega$. In this theorem $u(x,y)$ is considered to define a saddle surface if it is continuous and the function $u(x,y)-(ax+by+c)$ satisfies the weak maximum--minimum principle for all constants $a$, $b$, $c$. In particular, this will be the case if $u(x,y)$ satisfies (12.42) or represents a surface having non-positive Gaussian curvature. An elementary (but still not simple) proof is due to von Neumann [NU]. Hartman and Nirenberg [HN], [NI 4] have extended the result to higher dimensions.


% PDF page 329
\source{gilbarg-trudinger:page-0329}

With regard to the bounded slope condition, Hartman [HA] analyzes the relation between the regularity properties of $\varphi$ and $\partial\Omega$ when the function $\varphi$ satisfies a bounded slope condition over the boundary $\partial\Omega$ of a bounded convex domain $\Omega$ in $\R^n$. In the process he establishes the equivalence of the bounded slope and $(n+1)$-point conditions for $n\geq 2$; however, the relation between the constants in the two conditions remains unclear for $n>2$. Among the results we mention the following. (i) If $\partial\Omega\in C^{1,\alpha}$, $0\leq \alpha\leq 1$, and $\varphi$ satisfies a bounded slope condition over $\partial\Omega$, then $\varphi\in C^{1,\alpha}(\partial\Omega)$. (ii) If $\partial\Omega\in C^{1,1}$ and is uniformly convex, then $\varphi$ satisfies a bounded slope condition over $\partial\Omega$ if and only if $\varphi\in C^{1,1}$. The latter result follows from (i) and the fact that if $\Omega$ is any uniformly convex domain and $\varphi$ is the restriction to $\partial\Omega$ of a $C^{1,1}$ function, then $\varphi$ satisfies a bounded slope condition over $\partial\Omega$.

\vspace{1em}

\textbf{Problems}

\medskip

\noindent\textbf{12.1.} Show that the function $w(z)=z\log|z|$ is $(K,K')$-quasiconformal with $K=1$, $K'>0$ and does not satisfy (12.14) with the exponent $\alpha=1$.

\medskip

\noindent\textbf{12.2.} Prove Theorem 12.3 for quasiconformal mappings satisfying (12.2) in which $p,q$ are continuous and lie in $W^{1,2}_{\loc}$ (cf. [FS]).

\medskip

\noindent\textbf{12.3.} Let $Lu=au_{xx}+2bu_{xy}+cu_{yy}$ be uniformly elliptic in a domain $\Omega$ with $\gamma=\supp(\Lambda/\lambda)$, where $\lambda=\lambda(x,y)$ and $\Lambda=\Lambda(x,y)$ are the minimum and maximum eigenvalues of the coefficient matrix. If $Lu=0$, show that $p-iq=u_x-iu_y$ is $K$-quasiconformal in $\Omega$ for all $K\geq \frac12(\gamma+1/\gamma)$. [It suffices to show that

\[
\supp \frac{r^2+2s^2+t^2}{s^2-rt}=\frac{\lambda}{\Lambda}+\frac{\Lambda}{\lambda},
\]

where the supremum is taken over all $r,s,t$ satisfying $ar+2bs+ct=0$ for fixed $(x,y)\in\Omega$.] Hence infer that when $f=0$ Theorem 12.4 holds for any Hölder exponent $\alpha\leq 1/\gamma$ (cf. [TA] for $f\neq 0$).

\medskip

\noindent\textbf{12.4.} In Theorem 12.4, assume $f/\lambda\in L^p(\Omega)$, $p\geq 2$. If $\Omega'\Subset\Omega$, prove that for some $\alpha\in(0,1)$ we have

\[
|u|_{1,\alpha;\Omega'}\leq C\bigl(|u|_{0;\Omega}+\|f/\lambda\|_{L^p(\Omega)}\bigr),
\]

where $C=C(\gamma,p,\dist(\Omega',\partial\Omega))$ and $\alpha=\alpha(\gamma,p)$. (Either $u\in C^2(\Omega)$ or $u\in W^{2,2}(\Omega)$ may be assumed.)

\medskip

\noindent\textbf{12.5.} In Theorem 12.5, assume $\Omega\in C^{2,\beta}$, $\varphi\in C^{2,\beta}(\overline{\Omega})$ and $a,b,c,f\in C^{\beta}(\overline{\Omega}\times\R\times\R^2)$. By applying the $C^{1,\alpha}(\overline{\Omega})$ estimates discussed after Theorem 12.4, modify the proof of Theorem 12.5 to show that problem (12.29) has a solution in $C^{2,\beta}(\overline{\Omega})$. [Replace

% PDF page 330
\source{gilbarg-trudinger:page-0330}

C$^1_*$ and $C^{1,\alpha}_*$ by $C^{1,\alpha}(\overline{\Omega})$ and define the set $\mathcal{G}$ and the operator $T$ accordingly.
Use Theorem 6.14 in place of Theorem 6.13, and global in place of interior
Schauder estimates. A simpler proof is obtained by considering the family of
equations, $u=\sigma Tu$, $\sigma\in[0,1]$, corresponding to the Dirichlet problems

\[
Q_\sigma u = au_{xx}+2bu_{xy}+cu_{yy}+\sigma f = 0 \text{ in } \Omega,\qquad u=\sigma\varphi \text{ on } \partial\Omega,
\]

where $a=a(x,y,u,u_x,u_y)$, etc. Show that the solutions are uniformly bounded in
$C^{1,\alpha}(\overline{\Omega})$ and apply Theorem 11.3].

\textbf{12.6.}\ \ Assume the operator $Q$ in equation (12.40) is uniformly elliptic on bounded
subsets of $\Omega\times\mathbb{R}\times\mathbb{R}^2$. Prove Theorem 12.7 without the use of Theorem 12.5
by replacing the set $\mathcal{G}$ in Theorem 12.5 with

\[
\mathcal{G}'=\{v\in C^{1,\alpha}_* \mid |v|^{*}_{1,\alpha}\leq K',\ \ |Dv|\leq K\}
\]

where $K$ is the constant of the bounded slope condition and $K'$ is a constant
determined as in (12.32).

% PDF page 331
\source{gilbarg-trudinger:page-0331}

Chapter 13

\section*{H\"older Estimates for the Gradient}

In this chapter we derive interior and global H\"older estimates for the derivatives
of solutions of quasilinear elliptic equations of the form

\[
(13.1)\qquad Qu = a^{ij}(x,u,Du)D_{ij}u + b(x,u,Du) = 0
\]

in a bounded domain $\Omega$. From the global results we shall see that Step IV of the
existence procedure described in Chapter 11 can be carried out if, in addition
to the hypotheses of Theorem 11.4, we assume that either the coefficients $a^{ij}$ are in
$C^1(\overline{\Omega}\times\mathbb{R}\times\mathbb{R}^n)$ or that $Q$ is of divergence form or that $n=2$. The estimates of this
chapter will be established through a reduction to the results of Chapter 8, in
particular to Theorems 8.18, 8.24, 8.26 and 8.29.

\subsection*{13.1. \ Equations of Divergence Form}

Let us suppose now that $Q$ is equivalent to an elliptic operator of the form,

\[
(13.2)\qquad Qu = \operatorname{div} A(x,u,Du) + B(x,u,Du),
\]

where the vector function $A \in C^1(\Omega\times\mathbb{R}\times\mathbb{R}^n)$ and
$B \in C^0(\Omega\times\mathbb{R}\times\mathbb{R}^n)$. Then if $u \in C^1(\Omega)$ satisfies $Qu=0$ in $\Omega$, we have

\[
\int_{\Omega} \{A(x,u,Du)\cdot D\zeta - B(x,u,Du)\zeta\}\,dx = 0 \qquad \forall \zeta \in C_0^1(\Omega).
\]

Fixing $k$, $1\leq k\leq n$, replacing $\zeta$ by $D_k\zeta$, and integrating by parts, we then obtain

\[
\int_{\Omega} \{(D_{p_j}A^iD_{jk}u + \delta_kA^i)D_i\zeta + BD_k\zeta\}\,dx = 0 \qquad \forall \zeta \in C_0^1(\Omega),
\]

where $\delta_k$ is the differential operator defined by

\[
\delta_kA^i(x,z,p) = p_kD_zA^i(x,z,p) + D_{x_k}A^i(x,z,p)
\]
